\documentclass{article}
\usepackage[margin=1in]{geometry}

\title{External Review And Reproduction Certificates From Finite-Capacity Causal Networks}
\author{Abraham Greenman}
\date{2026}

\begin{document}
\maketitle

\begin{abstract}
This paper closes the Paper 16 conditional theorem for external review and
reproduction certificates downstream of the closed Paper 15
prediction-and-falsification protocol certificate. The result constructs a
finite, network-native certificate interface consisting of bounded review
records, bounded reproduction records, artifact and environment hash
descriptors, Paper 15 protocol references, stability and audit rows, and a
fail-closed no-hidden-claim audit. The closure is conditional and local: it
states that these certificate rows can be represented, checked, and assembled
without importing review acceptance, reproduction success, empirical adequacy,
physical validation, physical promotion, or a unified field theory claim.
\end{abstract}

\section{Scope}

Paper 16 begins only after the frozen Paper 15 conditional
prediction-and-falsification protocols certificate. Its task is not to obtain
external review acceptance or successful reproduction. Its task is to define
the finite certificate interface that can record review and reproduction
metadata while preserving the finite-capacity and non-promotion boundaries of
the upstream theorem ladder.

The upstream dependency is the frozen Paper 15 endpoint
\[
\texttt{Paper15PredictionFalsificationProtocolsTheoremContract.closed}
\]
with canonical proof
\[
\texttt{paper15\_pfp008\_final\_conditional\_certificate\_closes\_prediction\_falsification\_protocols\_theorem}.
\]

\section{Objects}

A Paper 16 certificate package is a finite tuple
\[
C=(U,R,D,H,K,S,A,F),
\]
where each component is a bounded certificate row or audit row:

\begin{itemize}
\item \(U\): upstream binding through the frozen Paper 15 certificate.
\item \(R\): finite external review and reproduction certificate record.
\item \(D\): finite reviewer, protocol, and provenance descriptors.
\item \(H\): finite artifact, environment, and hash descriptors.
\item \(K\): finite Paper 15 protocol compatibility reference.
\item \(S\): finite stability, auditability, and re-checkability behavior.
\item \(A\): no-hidden-claim import audit.
\item \(F\): final conditional certificate assembly.
\end{itemize}

All identifiers, reviewer labels, protocol labels, artifact labels,
environment descriptors, provenance labels, audit labels, status descriptors,
and hash descriptors are finite and bounded. Hash descriptors are integrity
bindings for audit rows only. Paper 15 compatibility is a reference relation
only. Reproducibility behavior means finite re-checkability of certificate
rows, not successful external reproduction.

\section{Closed Rungs}

\paragraph{ERRC-001.}
The upstream binding rung closes by recording the frozen Paper 1--15 chain,
the Paper 15 endpoint, and the Paper 15 final certificate while preserving
finite capacity, locality, bounded transfer, and all non-promotion boundaries.

\paragraph{ERRC-002.}
The finite certificate record rung closes by defining bounded certificate
identifiers, reviewer labels, protocol labels, artifact labels, environment
descriptors, reproduction-status descriptors, and audit-status descriptors.
Every row carries explicit no-recovery, no-acceptance, no-success,
no-validation, and no-promotion guards.

\paragraph{ERRC-003.}
The reviewer, protocol, and provenance descriptor rung closes by defining
bounded reviewer labels, reviewer-role labels, protocol labels, protocol-scope
labels, provenance source labels, provenance timestamp labels, and provenance
custodian labels. These descriptors are label and audit metadata only.

\paragraph{ERRC-004.}
The artifact, environment, and hash rung closes by defining bounded artifact
labels, artifact hash descriptors, environment descriptors, environment hash
descriptors, and non-success reproduction-status descriptors. Hashes bind
finite audit rows and do not state that reproduction succeeded.

\paragraph{ERRC-005.}
The Paper 15 compatibility rung closes by defining a bounded reference from
Paper 16 certificate rows to the frozen Paper 15 endpoint and final
certificate. The relation does not recover Paper 15 protocols and does not
claim prediction success, falsification success, review acceptance, or
reproduction success.

\paragraph{ERRC-006.}
The stability and auditability rung closes by defining finite audit snapshots
that preserve certificate identifiers, reviewer/protocol/provenance
descriptors, artifact/environment hash bindings, and Paper 15 compatibility
references. Re-checkability is finite row re-evaluation, not external
reproduction success.

\paragraph{ERRC-007.}
The no-hidden-claim audit closes by rejecting certificate recovery, protocol
recovery, review acceptance, reproduction success, benchmark success,
prediction success, falsification success, physical promotion, physical
validation, empirical adequacy, observed-catalog recovery, physical Standard
Model recovery, physical particle-excitation recovery, matter-field recovery,
gauge-field recovery, physical quantum dynamics, continuum quantum field
theory, simulation-only promotion, fit-only calibration, physical nature
realization, and unified-field promotion.

\paragraph{ERRC-008.}
The final conditional certificate rung closes by assembling ERRC-001 through
ERRC-007 and preserving every claim-boundary false flag.

\section{Theorem}

\textbf{External Review and Reproduction Certificates Theorem.}
Assume the frozen Paper 15 conditional prediction-and-falsification protocols
certificate and the closed ERRC-001 through ERRC-007 contracts. Then there is
a finite Paper 16 certificate package \(C\) such that:

\begin{enumerate}
\item all eight ERRC rungs are closed;
\item the Paper 16 external review and reproduction certificates theorem is
closed conditionally;
\item the closure is bounded, auditable, and network-native;
\item no certificate recovery, protocol recovery, review acceptance,
reproduction success, benchmark success, prediction success, falsification
success, empirical adequacy, physical validation, physical promotion,
physical nature realization, or unified field theory claim is introduced.
\end{enumerate}

\section{Proof}

Construct \(C=(U,R,D,H,K,S,A,F)\) from the canonical closed contracts for
ERRC-001 through ERRC-008.

ERRC-001 supplies \(U\), the upstream binding and claim-boundary preservation
contract. ERRC-002 supplies \(R\), the finite certificate row schema. ERRC-003
supplies \(D\), the finite reviewer, protocol, and provenance descriptor
schema. ERRC-004 supplies \(H\), the finite artifact, environment, and hash
descriptor schema. ERRC-005 supplies \(K\), the Paper 15 reference-only
compatibility relation. ERRC-006 supplies \(S\), the finite stability,
auditability, and re-checkability behavior. ERRC-007 supplies \(A\), the
fail-closed no-hidden-claim import audit. ERRC-008 supplies \(F\), the final
assembly certificate.

By construction each component is finite and bounded. The only upstream
mathematical endpoint referenced by the local certificate package is the
frozen Paper 15 endpoint. The compatibility relation records that endpoint
and certificate as finite labels; it does not convert protocol compatibility
into protocol recovery, prediction success, or falsification success.

The audit row \(A\) verifies that no component of \(C\) imports review
acceptance, reproduction success, physical validation, empirical adequacy,
physical promotion, physical nature realization, or unified-field promotion.
Therefore the final assembly row \(F\) may set the local theorem-closed flag
while all claim-boundary false flags remain false.

The Lean formal endpoint is:
\[
\texttt{paper16\_errc008\_final\_conditional\_certificate\_closes\_external\_review\_reproduction\_certificates\_theorem}.
\]
The Rust audit surface constructs the same finite certificate and verifies
that \(\texttt{closes\_paper16\_theorem()}\) returns true only when all rung
closures and non-promotion guards are present.

\section{Claim Boundary}

The theorem does not prove certificate recovery, review acceptance,
reproduction success, protocol recovery, benchmark success, prediction
success, falsification success, physical promotion, physical validation,
empirical adequacy, observed particle catalog recovery, physical Standard
Model content, physical particle excitations, physical matter fields,
physical gauge fields, physical quantum dynamics, continuum quantum field
theory, simulation-only promotion, fit-only calibration, physical nature
realization, or a unified field theory.

\section{Reproducibility Artifacts}

The formal proof lives in:
\[
\texttt{GPD/formal/FiniteCapacity/ExternalReviewReproductionCertificates.lean}.
\]
The executable audit guards live in:
\[
\texttt{rust/cclab\_accel/src/lib.rs}
\]
and
\[
\texttt{rust/cclab\_accel/tests/external\_review\_reproduction\_certificates.rs}.
\]
The required repository checks are:
\[
\texttt{make test-fast}
\qquad
\texttt{make lean-build}.
\]

\section{Conclusion}

Paper 16 is closed conditionally. It constructs a finite external review and
reproduction certificate interface downstream of Paper 15 while preserving
the theorem ladder's non-promotion boundary. The closure is an interface and
audit result only; it is not review acceptance, successful reproduction,
physical validation, or a unified field theory claim.

\end{document}
